\documentclass[a4paper,11pt]{article}
\usepackage{exam-style-842-unit}
\examinehead{842 信息化概论 · 第09单元练习册}{练习三 参考答案与解析}
\begin{document}

\unitanswerscover{842 信息化概论 · 单元练习册}{第09单元参考答案}{练习三（农业信息化现状、数字乡村与综合论述）的答案与解析；每道判断题均给出改正后的正确表述；论述题给出完整答题框架、给分点与分值分配、常见失分。}

\answerpart{一、单项选择题（10 题，每题 3 分，共 30 分）}

\q B。\begin{analysis}农业信息化是在农业全过程和农村经济社会各领域广泛应用信息技术、开发利用农业信息资源、发展农业信息服务的过程，覆盖生产、经营、管理与服务。A 只涉及生产环节的一项技术；C 只讲基础设施；D 只讲流通环节，三者都以偏概全。\end{analysis}

\q B。\begin{analysis}“最后一公里”指信息、技术与服务难以真正到达农户和田间地头：平台建在县里、数据停在部门，农户用不上、用不起、用不好。A 是农产品物流问题；C 是道路基础设施问题；D 是设备成本问题的一个方面，均不是“最后一公里”的本义。\end{analysis}

\q A。\begin{analysis}2019 年中共中央办公厅、国务院办公厅印发《数字乡村发展战略纲要》，明确提出数字乡村是乡村振兴的战略方向、建设数字中国的重要内容。B 是 2006 年国家层面的信息化战略文件；C 是 2021 年施行的法律；D 是“互联网+”行动的文件，均不是数字乡村建设的纲领性文件。\end{analysis}

\q B。\begin{analysis}数字乡村是伴随网络化、信息化和数字化在农业农村经济社会发展中的应用，以及农民现代信息技能提高而内生的农业农村现代化发展和转型进程，核心是数字技术赋能乡村全面振兴。A 只讲网络接入；C、D 只讲某一应用场景，均窄化了数字乡村的内涵。\end{analysis}

\q A。\begin{analysis}农业信息服务的内容包括政策法规、生产技术、市场行情与供求、气象与灾害预警、生产资料与金融保险、培训与咨询等。B、C、D 都只列了其中一项，不能涵盖农业信息服务的全部内容。\end{analysis}

\q A。\begin{analysis}信息进村入户工程按照“有场所、有人员、有设备、有宽带、有网页、有持续运营能力”的标准在村级建设信息服务站，统一称为益农信息社，提供公益服务、便民服务、电子商务和培训体验服务。B 农技推广站是农技推广机构；C 供销社是供销合作组织；D 农村信用社是农村金融机构，均不是信息进村入户的村级站点。\end{analysis}

\q B。\begin{analysis}12316 是全国农业系统公益服务统一专用号码，农民拨打后可咨询政策法规、生产技术、市场行情等问题。A 12315 是市场监管消费投诉举报热线；C 12345 是政务服务便民热线；D 12306 是铁路客票服务号码。\end{analysis}

\q A。\begin{analysis}区块链去中心化、不可篡改、可追溯、公开透明的特性，恰好契合农产品“从田间到餐桌”全过程信息记录与责任追溯的需求，因此农产品质量安全溯源是其最典型的农业应用。B 农田灌溉控制依靠物联网与自动控制；C 农机调度依靠定位与物联网；D 气象预报依靠气象监测与数值模型。\end{analysis}

\q A。\begin{analysis}农业物联网沿用物联网“感知层—网络层—支撑层—应用层”的体系架构（教材 7.1.2，印刷页 114—115）：感知层采集数据，网络层传输数据，支撑层提供计算与存储，应用层面向农业场景提供应用服务。B、C、D 是计算机网络协议的分层概念，不是物联网的组成层次。\end{analysis}

\q B。\begin{analysis}农业信息化是涉及基础设施、信息资源、人才、标准、资金、服务等多要素的系统工程，必须政府主导、多方参与、协同推进。A 单纯增加硬件会造成“重建设轻应用”；C 完全依靠市场难以解决公益性和区域不平衡问题；D 照搬国外模式不符合我国小农户为主的基本国情。\end{analysis}

\answerpart{二、多项选择题（3 题，每题 4 分，共 12 分）}

\q ABC。\begin{analysis}我国农业信息化的主要问题包括：农村信息基础设施仍较薄弱、存在“最后一公里”；农业信息资源分散、共享不足，数据孤岛突出；农业信息化复合型人才缺乏，农民数字技能不足；标准规范与资金投入不足等。D 错误：农业机械化水平过高不是农业信息化的问题，我国不少地区机械化、信息化水平仍有待提高。\end{analysis}

\q ABCD。\begin{analysis}数字乡村建设涵盖数字基础设施建设、智慧农业、乡村数字治理、信息惠民服务、乡村网络文化、数字生态等，目标是推动乡村生产、生活、治理、生态全面数字化。四项全对。\end{analysis}

\q ABC。\begin{analysis}区块链可用于农产品质量安全溯源（全链条记录、不可篡改）、农业供应链金融（确权与信用传递）、农村土地确权登记（权属信息可信存证）等。D 错误：农田自动灌溉控制依靠传感器、控制器与物联网，不属于区块链的典型应用。\end{analysis}

\answerpart{三、判断题（6 题，每题 2 分，共 12 分）}

\q ×。改正：农业信息化覆盖农业生产、经营、管理和服务全过程，并延伸到农村经济社会各领域，不只是生产环节的信息技术应用。\begin{analysis}把农业信息化等同于“生产环节信息化”，忽视了经营、管理、服务与乡村治理等环节，是概念窄化。\end{analysis}

\q ×。改正：数字乡村是以数字技术赋能乡村生产、生活、治理和生态，推动乡村全面振兴的系统工程，网络全覆盖只是其中的基础条件之一，还包括智慧农业、数字治理、信息惠民等内容。\begin{analysis}数字乡村不等于“通网”，还包括数据资源、应用场景、治理能力和农民数字素养等多方面建设。\end{analysis}

\q √。\begin{analysis}区块链的不可篡改、可追溯特性使农产品生产、加工、流通各环节信息难以被篡改，消费者和监管部门可追溯来源与责任，因此广泛用于农产品质量安全溯源。\end{analysis}

\q √。\begin{analysis}益农信息社是信息进村入户工程在村级设立的服务站点，承担公益服务、便民服务、电子商务和培训体验服务，是信息服务下沉到村的重要载体。\end{analysis}

\q ×。改正：农业信息服务不仅提供生产技术信息，还包括政策法规、市场行情、气象与灾害预警、生产资料与金融保险、培训与咨询等内容。\begin{analysis}只强调生产技术，忽视了市场、政策与风险信息，不能满足农户生产经营决策的完整需要。\end{analysis}

\q √。\begin{analysis}《数字乡村发展战略纲要》明确指出，数字乡村既是乡村振兴的战略方向，也是建设数字中国的重要内容。\end{analysis}

\answerpart{四、填空题（6 空，每空 1 分，共 6 分）}

\q （1）数字乡村。\begin{analysis}2019 年中共中央办公厅、国务院办公厅印发《数字乡村发展战略纲要》，是数字乡村建设的纲领性文件。\end{analysis}

\q （2）12316。\begin{analysis}12316 是全国农业系统公益服务统一专用号码，为农民提供政策、技术、市场等咨询服务。\end{analysis}

\q （3）益农信息社。\begin{analysis}信息进村入户工程在村级设立的服务站点统一称为益农信息社。\end{analysis}

\q （4）应用层。\begin{analysis}农业物联网由感知层、网络层、支撑层和应用层组成（教材 7.1.2）。\end{analysis}

\q （5）可追溯。\begin{analysis}区块链具有去中心化、不可篡改、可追溯、公开透明等特点，因而适合农产品溯源。\end{analysis}

\q （6）气象预警（或气象与灾害预警）。\begin{analysis}农业信息服务内容包括政策法规、生产技术、市场行情、气象与灾害预警、金融保险和培训咨询等。\end{analysis}

\answerpart{五、简答题（3 题，每题 8 分，共 24 分）}

\q 参考答案要点：
（1）现状（3 分）：农村信息基础设施明显改善，行政村宽带和移动网络覆盖率大幅提升；农业农村大数据平台、益农信息社、12316 热线等服务体系初步建立；农业物联网、智能农机、无人农场在示范基地加快应用；农产品电商和数字乡村建设深入推进，政策支持力度不断加大。
（2）主要问题（5 分，答出 5 点即可）：一是农村信息基础设施仍较薄弱，存在“最后一公里”问题；二是农业信息资源分散、标准不统一、共享不足，数据孤岛突出；三是农业信息化复合型人才缺乏，农民数字技能不足；四是建设与运维资金投入不足，投融资机制不健全；五是信息技术应用与农业生产实际结合不紧，重建设轻应用；六是区域发展不平衡，小农户应用门槛高，服务体系和商业模式不健全。

\q 参考答案要点（答出 4～5 点即可，每点 2 分，总分不超过 8 分）：
（1）政府主导、多元投入：加大财政投入，引导社会资本参与，建立可持续的建设运维机制；
（2）补齐基础设施短板：推进农村宽带、5G、数据中心等新基建，解决“最后一公里”；
（3）整合信息资源、健全标准：建设农业农村大数据平台，统一数据标准与接口，打破数据孤岛，推动共享开放；
（4）培育人才队伍：加强基层农技人员信息化培训，培养复合型人才，提升农民数字素养；
（5）示范推广、服务下沉：建设示范基地和益农信息社，发展农业社会化服务，降低小农户应用门槛；
（6）推动技术与生产融合：面向水稻、柑橘、生猪等主导产业开发实用应用，形成可持续商业模式；
（7）完善法规与安全保障：加强数据安全、个人信息保护和涉农平台监管。

\q 参考答案要点：
（1）内涵（4 分）：数字乡村是伴随网络化、信息化和数字化在农业农村经济社会发展中的应用，以及农民现代信息技能提高而内生的农业农村现代化发展和转型进程；既是乡村振兴的战略方向，也是建设数字中国的重要内容。其内涵包括数字基础设施、智慧农业、乡村数字治理、信息惠民服务、乡村网络文化和数字生态等方面。
（2）实现途径（4 分，答出 4 点即可）：一是加强农村数字基础设施建设（宽带、5G、数据中心、物联网感知设施）；二是发展智慧农业，建设农业农村大数据平台，推广物联网、智能农机和精准作业；三是推进乡村数字治理，开展数字党建、网格化管理、政务服务“一网通办”和应急管理数字化；四是深化信息惠民，建设益农信息社、发展农产品电商和“互联网+医疗、教育、文化”服务；五是培育数字化人才、完善政策与资金保障、健全数据标准与安全制度。

\answerpart{六、论述题（1 题，共 16 分）}

\q 参考答案与答题框架（按“破题—背景—技术—应用—问题—对策—升华”展开）：

\textbf{【破题句】}农业信息化是乡村振兴的重要支撑和农业强国建设的重要抓手。物联网、大数据、人工智能、区块链等信息技术正从生产端和治理端同时重塑农业与乡村，但“最后一公里”、数据孤岛、人才与资金短板仍制约其效能，必须技术、应用、问题、对策贯通作答。

\textbf{一、背景（2 分）}
\begin{analysis}
关键词：乡村振兴战略“二十字”总要求、数字乡村发展战略纲要、小农户与现代农业衔接、农业强国。
要点：乡村振兴要求产业兴旺、治理有效，数字乡村为农业农村现代化提供新动能；我国以小农户为主的经营格局决定了信息化必须走“技术下沉、服务到户”的道路。
\end{analysis}

\textbf{二、四类信息技术及其作用（4 分）}
\begin{analysis}
每类 1 分，须点明技术特征与解决的问题：
（1）物联网：全面感知、可靠传输、智能控制，解决农业生产“看不见、控不准”的问题；
（2）大数据：汇聚多源数据、分析预测，解决“数据散、算不清”的问题；
（3）人工智能：图像识别、预测与决策，解决“靠经验、算不准”的问题；
（4）区块链：去中心化、不可篡改、可追溯，解决“信不过、追不到”的问题。
\end{analysis}

\textbf{三、应用：生产端与治理端（4 分）}
\begin{analysis}
生产端（2 分，须结合具体产业举例）：
（1）水稻：育秧大棚温湿度物联网监控、智能灌溉、无人机植保与遥感长势监测；
（2）柑橘：AI 图像识别病虫害、按处方变量施药、果园水肥一体化；
（3）生猪：养殖环境自动控制、个体健康监测、区块链耳标溯源；
（4）茶叶：区块链记录种植、加工、检测全过程，支撑品牌与电商销售。
治理端（2 分）：
（1）数字党建、网格化管理、视频监控“雪亮工程”、应急指挥数字化；
（2）“互联网+政务服务”向乡村延伸，益农信息社提供公益、便民、电商与培训服务；
（3）农村集体资产、土地确权等数据上链存证，提升治理透明度。
\end{analysis}

\textbf{四、问题（3 分）}
\begin{analysis}
答出 3 点、每点 1 分：
（1）基础设施与“最后一公里”：网络、电力、感知设施不完善，服务难到户；
（2）数据与标准：资源分散、标准不统一、共享机制缺失，数据孤岛突出；
（3）人才与资金：复合型人才缺乏、农民数字技能不足，建设运维资金与商业模式不健全；
（4）技术与生产脱节：重建设轻应用，小农户经营规模小、应用成本高；
（5）区域不平衡与安全风险：城乡、区域差距明显，数据安全与个人信息保护有待加强。
\end{analysis}

\textbf{五、对策（2 分）}
\begin{analysis}
答出 2～3 点、与问题对应：
（1）政府主导、多元投入，补齐新基建设施短板，建立长效运维机制；
（2）建设农业农村大数据平台，统一标准、打破孤岛，推动数据共享开放；
（3）培育人才、提升农民数字素养，发展社会化服务，降低小农户应用门槛；
（4）以产业需求为导向开展示范推广，形成“政府引导、市场运作、农民受益”的可持续发展模式。
\end{analysis}

\textbf{六、升华（1 分）}
\begin{analysis}
关键词：以数据为关键要素、以农民为主体、因地制宜、走中国特色智慧农业道路。
收束句：农业信息化不是技术的简单叠加，而是以数据为要素、以产业为根基、以农民为主体的系统性变革；只有把技术嵌入生产实际、把服务送到田间地头，才能真正赋能乡村振兴。
\end{analysis}

\textbf{【给分点与分值分配】}背景 2 分；四类技术 4 分（每类 1 分）；应用 4 分（生产端 2 分、治理端 2 分）；问题 3 分；对策 2 分；升华 1 分，合计 16 分。

\textbf{【常见失分】}一是只堆“物联网、大数据、人工智能、区块链”名词，不说明各自解决什么问题，技术介绍与农业应用“两张皮”；二是通篇只写技术应用，没有问题与对策，或问题与对策不能一一对应；三是不结合乡村振兴、数字乡村背景，把论述写成技术说明文；四是举例笼统（“某地某农场”），没有落到水稻、柑橘、生猪、茶叶等具体产业与环节；五是对策空喊“加强重视、加大投入”，缺少主体、抓手和机制；六是结构混乱、不分点、字数不足，缺少破题与收束。

\end{document}
